% Source: contrastif_learning.tex
\begin{notechapter}
\chapter{Contrastive Learning: An Overview}\label{ch:contrastive-overview}

\section{What is Contrastive Learning?}
Contrastive learning is a self-supervised or unsupervised learning 
technique used primarily in machine learning, particularly in 
representation learning. The goal of contrastive learning is to learn 
meaningful representations of data by contrasting similar (positive) 
pairs against dissimilar (negative) pairs. It has gained significant 
attention in recent years, especially in computer vision and natural 
language processing, for its ability to learn high-quality 
representations without requiring labeled data.

\section{Key Idea of Contrastive Learning}
The core idea is to pull similar data points (positive pairs) closer 
together in the representation space while pushing dissimilar data 
points (negative pairs) farther apart. This is achieved by optimizing a 
contrastive loss function.

\section{How Contrastive Learning Works}
\begin{itemize}
    \item \textbf{Positive and Negative Pairs}:
          \begin{itemize}
              \item \textbf{Positive Pair}: Two data points that are 
                    semantically similar (e.g., two augmented views of 
                    the same image or two sentences with similar 
                    meanings).
              \item \textbf{Negative Pair}: Two data points that are 
                    semantically dissimilar (e.g., two different images 
                    or unrelated sentences).
          \end{itemize}
    \item \textbf{Representation Learning}:
          \begin{itemize}
              \item A neural network (e.g., a convolutional neural 
                    network or transformer) is used to encode data 
                    points into a lower-dimensional representation 
                    space.
              \item The goal is to ensure that representations of 
                    positive pairs are close together, while 
                    representations of negative pairs are far apart.
          \end{itemize}
    \item \textbf{Contrastive Loss Function}:
          \begin{itemize}
              \item The most commonly used loss function is the 
                    \textbf{InfoNCE loss} (Noise Contrastive 
                    Estimation), which is defined as:
                    \[
                    \mathcal{L} = -\log \frac{\exp(\text{sim}(z_i, z_j) 
                    / \tau)}{\sum_{k=1}^N \exp(\text{sim}(z_i, z_k) / 
                    \tau)}
                    \]
                    Here:
                    \begin{itemize}
                        \item \(z_i\) and \(z_j\) are representations 
                              of a positive pair.
                        \item \(z_k\) are representations of negative 
                              pairs.
                        \item \(\text{sim}(\cdot)\) is a similarity 
                              function (e.g., cosine similarity).
                        \item \(\tau\) is a temperature parameter that 
                              controls the sharpness of the 
                              distribution.
                    \end{itemize}
          \end{itemize}
    \item \textbf{Data Augmentation}:
          \begin{itemize}
              \item In computer vision, contrastive learning often 
                    relies on data augmentation to create positive 
                    pairs. For example, two augmented views of the same 
                    image (e.g., cropped, rotated, or color-jittered 
                    versions) are treated as a positive pair.
              \item In NLP, positive pairs can be created using 
                    paraphrasing or sentence transformations.
          \end{itemize}
\end{itemize}

\section{Applications of Contrastive Learning}
\begin{itemize}
    \item \textbf{Computer Vision}:
          \begin{itemize}
              \item \textbf{SimCLR}: A popular framework for 
                    contrastive learning in vision, which uses data 
                    augmentation to create positive pairs.
              \item \textbf{MoCo (Momentum Contrast)}: Another framework 
                    that uses a momentum encoder to improve the quality 
                    of negative samples.
              \item \textbf{BYOL (Bootstrap Your Own Latent)}: A variant 
                    that does not require negative samples.
          \end{itemize}
    \item \textbf{Natural Language Processing (NLP)}:
          \begin{itemize}
              \item \textbf{Sentence Embeddings}: Contrastive learning 
                    is used to learn sentence representations by 
                    contrasting semantically similar sentences.
              \item \textbf{Cross-Modal Learning}: Contrastive learning 
                    is used to align representations across modalities 
                    (e.g., text and images).
          \end{itemize}
    \item \textbf{Other Domains}:
          \begin{itemize}
              \item \textbf{Audio}: Learning representations of audio 
                    signals by contrasting similar audio clips.
              \item \textbf{Graphs}: Learning node or graph 
                    representations by contrasting similar graph 
                    structures.
          \end{itemize}
\end{itemize}

\section{Advantages of Contrastive Learning}
\begin{itemize}
    \item \textbf{Self-Supervised Learning}: It does not require labeled 
          data, making it highly scalable.
    \item \textbf{High-Quality Representations}: It learns 
          representations that capture semantic similarities and are 
          useful for downstream tasks.
    \item \textbf{Flexibility}: It can be applied to various data types, 
          including images, text, audio, and graphs.
\end{itemize}

\section{Challenges in Contrastive Learning}
\begin{itemize}
    \item \textbf{Choice of Positive and Negative Pairs}:
          \begin{itemize}
              \item The quality of learned representations depends 
                    heavily on how positive and negative pairs are 
                    defined.
          \end{itemize}
    \item \textbf{Scalability}:
          \begin{itemize}
              \item The number of negative samples can grow very large, 
                    making computation expensive.
          \end{itemize}
    \item \textbf{Bias in Negative Sampling}:
          \begin{itemize}
              \item If negative samples are not diverse enough, the 
                    model may learn suboptimal representations.
          \end{itemize}
\end{itemize}

\section{Summary}
Contrastive learning is a powerful technique for learning meaningful 
representations of data by contrasting similar and dissimilar pairs. It 
has been successfully applied in various domains, including computer 
vision and NLP, and is a key component of many state-of-the-art 
self-supervised learning methods.
\end{notechapter}
